\section{Introduction}

Agent runtimes increasingly retain an append-only record of model calls, tool
results, human decisions, state updates, and coordination events.  Event
sourcing makes several desirable properties seem immediate.  If state is a
fold over the log, a run can be replayed.  If a child log shares a prefix with a
parent, a run can be forked without recomputing that prefix.  If causes are
recorded, artifacts can be traced back to model and tool activity.  ActiveGraph
states these properties as consequences of making the log the source of truth
\cite{nakajima2026log}.

This paper formalizes and empirically checks the replay and fork properties
claimed on architectural grounds in Nakajima (2026).  The relationship is
deliberate: the earlier work describes a deployed architecture, while the
present work asks which side conditions make its architectural claims sound.
Negative results about the originating runtime are therefore part of the
contribution rather than exceptions to be hidden.

Two gaps motivate the work.  First, ``deterministic replay'' often means
deterministic replay \emph{after oracle outputs have been fixed}.  This is
useful, but the phrase suppresses operational questions.  Was the model output
captured faithfully?  Is a request retained without its result?  Does replay
serve the recorded result or make a new call?  Is the runtime itself
schedule-independent, or does a particular FIFO order merely produce one
repeatable execution?

Second, event sourcing in a business application normally reconstructs
application state; it does not claim to reverse the external world.  That
distinction becomes load-bearing for agents.  A retrospectively discarded log
suffix may contain an email, a charge, a delete, a publication, or a billable
model call.  The forked trace can omit the effect, but the world cannot be
assumed to be one in which the effect never happened.  A correct fork contract
must say which property it preserves.

The initial version of this project asked whether

\begin{equation}
  A : S \times E \rightarrow S \times F^*
\end{equation}

was the irreducible agent primitive.  That framing is not novel.  The shape is
a Mealy-style transition system \cite{mealy1955method}, closely resembles the
Elm update architecture \cite{elmarchitecture}, and appears in functional
event sourcing through the Decider pattern \cite{chassaing2021decider}.  We
concede the minimality lane and instead exploit the effects/events seam that
becomes visible when agents interact with irreversible external systems.

\paragraph{Contributions.}
This draft makes four contributions.

\begin{enumerate}[leftmargin=*]
  \item It gives an executable operational model in which scheduler order is
  explicit.  The model separates termination, blocking on an external effect,
  and bounded divergence.

  \item It preserves counterexamples showing that general settlement is neither
  terminating nor confluent.  It further shows that pairwise-disjoint
  \emph{declared writes} are insufficient when behavior reads and emitted
  triggers interfere.  A restricted isolated-writer family does converge, but
  we do not claim a minimal general condition.

  \item It defines property-relative fork assessment over three orthogonal
  effect dimensions: footprint, replay source, and lifecycle.  This separates
  exact retained-prefix projection from execution replay, continuation, and a
  counterfactual-world claim.

  \item It checks the laws against 5,564 runtime logs, 314,348 source events,
  every possible cut in each log, and 121 observed forks.  The validation finds
  both a strong positive result---zero structural mismatches in the observed
  shared prefixes---and sharp limitations on which hypothetical cuts license
  continuation or counterfactual claims.
\end{enumerate}

\paragraph{Genre.}
This is not a mechanized-proof paper.  Laws are stated as propositions, checked
over generated families with FsCheck, and evaluated exhaustively over the cuts
of the available logs.  A failed law produces a minimized, named regression
fixture.  Passing checks are finite evidence over the specified generators and
artifacts, not universal proofs.  This choice is intentional: calculi with
machine-checked composition, termination, and information-flow results already
exist \cite{liu2026lambdaa,garby2026llmbda}.  Our comparative advantage is an
executable runtime contract connected to real runtime artifacts.

\paragraph{Roadmap.}
Section~\ref{sec:model} defines the executable model and observations.
Section~\ref{sec:confluence} reports termination and confluence results.
Section~\ref{sec:fork} defines fork properties and effect conditions.
Section~\ref{sec:grades} connects those conditions to replayability grades.
Sections~\ref{sec:method} and~\ref{sec:validation} describe the artifact and
empirical validation.  We then position the work, discuss implications, and
state limitations.
